\documentclass{scrartcl}
\usepackage{graphicx}  % required for inserting images

% for cyrillic and Bulgarian language
\usepackage[utf8]{inputenc}
\usepackage[T2A]{fontenc}
\usepackage[bulgarian]{babel}

% math packages
\usepackage{amsmath}
\usepackage{float}
\usepackage{amsfonts}
\usepackage{amssymb}
\usepackage{amsthm}
\usepackage{cancel}

\usepackage{ragged2e}  % adds FlushLeft

\usepackage{listings}
\usepackage{xcolor}

\definecolor{codegreen}{rgb}{0,0.6,0}
\definecolor{codegray}{rgb}{0.5,0.5,0.5}
\definecolor{codepurple}{rgb}{0.58,0,0.82}
\definecolor{backcolour}{rgb}{0.95,0.95,0.92}

\lstdefinestyle{mystyle}{
    backgroundcolor=\color{backcolour},   
    commentstyle=\color{codegreen},
    keywordstyle=\color{magenta},
    numberstyle=\tiny\color{codegray},
    stringstyle=\color{codepurple},
    basicstyle=\ttfamily\footnotesize,
    breakatwhitespace=false,         
    breaklines=true,                 
    captionpos=b,                    
    keepspaces=true,                 
    numbers=left,                    
    numbersep=5pt,                  
    showspaces=false,                
    showstringspaces=false,
    showtabs=false,                  
    tabsize=2
}
\lstset{style=mystyle}

\usepackage[colorlinks=true, linkcolor=blue, urlcolor=cyan]{hyperref}

\newtheorem{definition}{Дефиниция}
\newtheorem{theorem}{Теорема}

\title{Бази от данни - релационен модел на данните}
\author{Ивайло Андреев + chat gpt 5.5, DeepSeek-V3, gemini 3 Flash}
\date{25 май 2026}

\begin{document}

\maketitle

\tableofcontents
\newpage

\begin{FlushLeft}

\section{Въведение}

\subsection{База данни и СУБД}

\textbf{\textit{База данни (БД)}} е организирана колекция от свързани данни, предназначена за ефективно съхранение, търсене и управление в информационни системи.

\textbf{\textit{Система за управление на база данни (СУБД)}} е софтуерен инструмент, който осигурява средства за дефиниране, съхранение, извличане и манипулиране на данни в БД. Примери: MySQL, PostgreSQL, Oracle, SQL Server, MongoDB.

\subsection{Модели на бази данни}

Моделът на БД определя логическата структура на данните и начина на тяхната организация. Основните категории са:

\begin{itemize}
    \item \textbf{Релационен модел} --- Представя данните като таблици (релации) с редове (кортежи) и колони (атрибути). Манипулацията се извършва чрез SQL. Подходящ за данни със стабилна, строго дефинирана структура.
\end{itemize}

Нерелационните модели (NoSQL) съществуват като алтернатива, но в тази тема фокусът е върху релационния модел.

\section{Релационен модел на данните}

\subsection{Домейн}

\begin{definition}
\textbf{\textit{Домейн (Област)}} е множество от допустими стойности за дадена величина. Релационният модел изисква стойностите на всички елементи в кортежите да бъдат \textbf{атомарни} (неделими на части).
\end{definition}

Домейнът обикновено дефинира скаларния тип данни (INTEGER, VARCHAR, DATE и т.н.) и допълнителни ограничения.

\subsection{Релация}

\begin{definition}
\textbf{\textit{Релация}} е подмножество на декартово произведение $D_1 \times D_2 \times \cdots \times D_n$, където $D_i$ са домейни. В контекста на БД релацията представлява таблица с редове (кортежи) и колони (атрибути).
\end{definition}

\subsection{Кортеж}

\begin{definition}
\textbf{\textit{Кортеж}} е всеки елемент на релация, т.е. наредена $n$-торка от стойности. Тъй като релацията е множество, всеки кортеж трябва да е уникален.
\end{definition}

\textbf{\textit{Екземпляр на релация}} е множеството от всички текущо съхранени кортежи. \textbf{\textit{Кардиналност}} е броят на кортежите в екземпляра.

\subsection{Атрибути}

\begin{definition}
\textbf{\textit{Атрибут}} е наименование на колона в релация. Всеки атрибут $A_i$ е асоцииран с домейн $D_i$. Атрибутите описват характеристиките на съхранените обекти.
\end{definition}

\subsection{Схема на релация}

\begin{definition}
\textbf{\textit{Схема на релация}} $R$ се дефинира чрез име $R$, множество от атрибути $\{A_1, A_2, \ldots, A_n\}$ и съответните домейни. Обозначение: $R(A_1, A_2, \ldots, A_n)$.
\end{definition}

Свързани понятия:

\begin{itemize}
    \item \textbf{Степен (Arity)} --- Броят на атрибутите в схемата. Статична характеристика.
    \item \textbf{Кардиналност} --- Броят на кортежите в момента. Динамична характеристика.
\end{itemize}

\subsection{Схема на релационна база данни}

\begin{definition}
\textbf{\textit{Схема на релационна БД}} е съвкупност от всички схеми на релации, включително ограниченията за интегритет (например външни ключове), които управляват връзките между релациите.
\end{definition}

\textbf{\textit{Екземпляр на БД}} е съвкупност от екземплярите на всички релации в дадения момент.

\subsection{Проектиране на релационна БД}

Проектирането обикновено включва:

\begin{enumerate}
    \item Определяне какви данни да се съхраняват.
    \item Дефиниране на връзки между данните.
    \item Определяне на ограничения (ключове, домейни и т.н.).
\end{enumerate}

В практиката се използва \textbf{Entity/Relationship (E/R) модел}, който се преобразува към релационна схема на БД чрез:

\begin{itemize}
    \item Всяка същност (Entity) се превръща в релация със същите атрибути.
    \item Всяка връзка (Relationship) се превръща в релация, чиито атрибути са ключовете на свързаните същности.
\end{itemize}

\subsection{Ключове}

\begin{definition}
\textbf{\textit{Суперключ}} е множество от атрибути, което еднозначно определя всеки кортеж в релацията.
\end{definition}

\begin{definition}
\textbf{\textit{Кандидат-ключ}} е минимален суперключ.
\end{definition}

\begin{definition}
\textbf{\textit{Първичен ключ}} е избран кандидат-ключ, използван за идентификация на кортежите.
\end{definition}

\begin{definition}
\textbf{\textit{Външен ключ}} е атрибут или множество от атрибути в релация $R$, които реферират към първичен ключ в друга (или същата) релация.
\end{definition}

\subsection{Интегритетни ограничения и NULL}

\begin{itemize}
    \item \textbf{Entity Integrity} --- първичният ключ не може да има \texttt{NULL} стойност.
    \item \textbf{Referential Integrity} --- всяка ненулева стойност на външен ключ трябва да съответства на съществуващ ключ в реферираната релация.
\end{itemize}

\textbf{\texttt{NULL}} означава липсваща, неизвестна или неприложима стойност. Тя не е число, не е празен низ и не е нула. При сравнения с \texttt{NULL} се използват \texttt{IS NULL} / \texttt{IS NOT NULL}.

\section{Видове операции}

\subsection{DDL (Data Definition Language)}

DDL операциите служат за дефиниране на схемата на БД. Включват CREATE, DROP и ALTER.

\begin{lstlisting}[language=SQL, caption=DDL examples]
-- Create table Employees
CREATE TABLE Employees (
    employee_id INT PRIMARY KEY,
    first_name VARCHAR(50) NOT NULL,
    last_name VARCHAR(50) NOT NULL,
    salary DECIMAL(10,2),
    department_id INT
);

-- Add new column
ALTER TABLE Employees ADD hire_date DATE;

-- Drop column
ALTER TABLE Employees DROP COLUMN hire_date;

-- Drop table
DROP TABLE Employees;
\end{lstlisting}

\subsection{DML (Data Manipulation Language)}

DML операциите служат за манипулация на данните: SELECT, INSERT, UPDATE, DELETE.

\begin{lstlisting}[language=SQL, caption=DML examples]
-- Select rows
SELECT employee_id, first_name, last_name
FROM Employees
WHERE salary > 50000;

-- Insert row
INSERT INTO Employees (employee_id, first_name, last_name, salary)
VALUES (1, 'John', 'Doe', 75000);

-- Update rows
UPDATE Employees
SET salary = 80000
WHERE employee_id = 1;

-- Delete rows
DELETE FROM Employees
WHERE salary < 30000;
\end{lstlisting}

\section{Заявки към релационна БД}

\subsection{Основни примери}

\begin{lstlisting}[language=SQL, caption=SELECT with JOIN]
-- INNER JOIN
SELECT 
    e.first_name, 
    e.last_name, 
    d.department_name
FROM Employees e
INNER JOIN Departments d ON e.department_id = d.department_id
WHERE e.salary > 60000;

-- LEFT JOIN
SELECT 
    d.department_name,
    e.first_name,
    e.last_name
FROM Departments d
LEFT JOIN Employees e ON d.department_id = e.department_id
ORDER BY d.department_name;
\end{lstlisting}

\begin{lstlisting}[language=SQL, caption={INSERT, UPDATE, DELETE examples}]
INSERT INTO Employees (employee_id, first_name, last_name, salary, department_id)
VALUES (10, 'Anna', 'Petrova', 3200.00, 3);

UPDATE Employees
SET salary = 3500.00
WHERE employee_id = 10;

DELETE FROM Employees
WHERE employee_id = 10;
\end{lstlisting}

\section{Тригери}

\begin{definition}
\textbf{\textit{Тригер (Trigger)}} е специален тип съхранена процедура, която автоматично се активира като отговор на определено събитие (INSERT, UPDATE, DELETE) в дадена таблица.
\end{definition}

\subsection{Пример за тригер}

\begin{lstlisting}[language=SQL, caption=Trigger: log salary changes]
-- Log salary changes into Salary_Log
CREATE TRIGGER log_salary_change
AFTER UPDATE OF salary ON Employees
FOR EACH ROW
BEGIN
    INSERT INTO Salary_Log (employee_id, old_salary, new_salary, change_date)
    VALUES (NEW.employee_id, OLD.salary, NEW.salary, NOW());
END;
\end{lstlisting}

\section{Релационна алгебра}

\textbf{\textit{Релационна алгебра}} е математическа система, която дефинира операции над релации. Служи като основа на SQL заявките.

\textbf{Важно:} SQL е практически език, вдъхновен от релационната алгебра, но не е напълно еквивалентен на нея. Аналогиите $\pi \leftrightarrow$ \texttt{SELECT} и $\sigma \leftrightarrow$ \texttt{WHERE} са полезни, но не са пълна формална идентичност.

\subsection{Съвместими релации}

\begin{definition}
Релациите $R(A_1, \ldots, A_n)$ и $S(B_1, \ldots, B_m)$ са \textbf{\textit{съвместими}}, ако:
\begin{itemize}
    \item $n = m$ (еднакъв брой атрибути)
    \item $\text{dom}(A_i) = \text{dom}(B_i)$ за всеки $i$ (еднакви домейни)
    \item Редът на атрибутите съвпада
\end{itemize}
\end{definition}

\subsection{Основни операции}

Основните операции в Релационната алгебра образуват базис и всички останали операции могат да се изразят чрез базисните.

\subsubsection{Обединение (Union)}

\begin{definition}
\textbf{\textit{Обединение}} $R \cup S$ на съвместими релации $R$ и $S$ съдържа всички кортежи, които се намират в $R$ или $S$ (или в двете).
\end{definition}

Свойства: бинарна, комутативна, асоциативна.

\begin{lstlisting}[language=SQL, caption=UNION]
SELECT employee_id, name FROM full_time_employees
UNION
SELECT employee_id, name FROM part_time_employees;
\end{lstlisting}

\subsubsection{Разлика (Difference)}

\begin{definition}
\textbf{\textit{Разлика}} $R \setminus S$ съдържа всички кортежи на $R$, които не се намират в $S$.
\end{definition}

Свойства: бинарна, некомутативна.

\begin{lstlisting}[language=SQL, caption=EXCEPT or MINUS]
SELECT employee_id, name FROM all_employees
EXCEPT
SELECT employee_id, name FROM terminated_employees;
\end{lstlisting}

\subsubsection{Проекция (Projection)}

\begin{definition}
\textbf{\textit{Проекция}} $\pi_{A_{i_1}, \ldots, A_{i_k}}(R)$ на релация $R$ с атрибути $\{A_1, \ldots, A_n\}$ е релация, съдържаща само избраните атрибути $\{A_{i_1}, \ldots, A_{i_k}\}$.
\end{definition}

Еквивалентна е на SQL операцията SELECT с избрани колони.

Ако проекцията не съдържа ключови атрибути, в SQL могат да възникнат дублиращи се редове, затова за пълна еквивалентност с релационната алгебра често се използва DISTINCT.

\begin{lstlisting}[language=SQL, caption=Projection (SELECT)]
SELECT DISTINCT first_name, last_name
FROM Employees;
\end{lstlisting}

\subsubsection{Селекция (Selection)}

\begin{definition}
\textbf{\textit{Селекция}} $\sigma_P(R)$ е множество от всички кортежи в $R$, които удовлетворяват предиката $P$.
\end{definition}

Е еквивалентна на SQL WHERE клаузула.

\begin{lstlisting}[language=SQL, caption=Selection (WHERE)]
SELECT *
FROM Employees
WHERE salary > 50000 AND department_id = 2;
\end{lstlisting}

\subsubsection{Преименуване (Renaming)}

\begin{definition}
\textbf{\textit{Преименуване}} $\rho_{R'/R}(R)$ или в SQL AS позволяват смяна на имена на релации и атрибути без промяна на данните.
\end{definition}

\begin{lstlisting}[language=SQL, caption=Renaming]
SELECT e.employee_id AS emp_id, e.first_name AS name
FROM Employees AS e;
\end{lstlisting}

\subsubsection{Декартово произведение (Cartesian Product)}

\begin{definition}
\textbf{\textit{Декартово произведение}} $R \times S$ съдържа всички възможни комбинации на кортежи от $R$ с кортежи от $S$.
\end{definition}

Ако $R$ има $m$ редове и $S$ има $n$ редове, $R \times S$ има $m \times n$ редове.

В релационната алгебра атрибути с едно и също име в различни релации трябва формално да се преименуват преди декартовото произведение между тях. В SQL това е позволено, но при проекция трябва да се уточни към коя таблица принадлежи колоната.

\begin{lstlisting}[language=SQL, caption=Cartesian product]
SELECT e.employee_id, d.department_id
FROM Employees e, Departments d;
-- Or: FROM Employees e CROSS JOIN Departments d;
\end{lstlisting}

\subsection{Допълнителни операции}

\subsubsection{Сечение (Intersection)}

\begin{definition}
\textbf{\textit{Сечение}} $R \cap S$ съдържа само кортежите, които се намират едновременно в $R$ и $S$.
\end{definition}

Свойства: комутативна, асоциативна. Може да се изрази чрез разлика: $R \cap S = R \setminus (R \setminus S)$.

\begin{lstlisting}[language=SQL, caption=INTERSECT]
SELECT employee_id FROM employees_2023
INTERSECT
SELECT employee_id FROM employees_2024;
\end{lstlisting}

\subsubsection{Съединение (Join)}

\begin{definition}
\textbf{\textit{Theta Join}} $R \underset{P}{\bowtie} S$ е множество от всички кортежи от $R \times S$, които удовлетворяват предиката $P$.
\end{definition}

Може да се изрази чрез декартово произведение и селекция: $R \underset{P}{\bowtie} S = \sigma_P(R \times S)$.

\textbf{\textit{Естествено съединение}} $R \bowtie S$ е частен случай на theta join, където съединяваме релации по атрибути с еднакви имена. Резултатът съдържа само един екземпляр от общите атрибути.

\begin{lstlisting}[language=SQL, caption=JOIN operations]
-- Theta Join (INNER JOIN)
SELECT *
FROM Employees e, Departments d
WHERE e.department_id = d.department_id;

-- Equivalent form:
SELECT *
FROM Employees e
INNER JOIN Departments d ON e.department_id = d.department_id;

-- Natural join (same attribute names)
SELECT *
FROM Employees
JOIN Departments USING (department_id);
\end{lstlisting}

\subsubsection{Частно (Division)}

\begin{definition}
\textbf{\textit{Частно}} $R \div S$ (където атрибутите на $S$ са подмножество от атрибутите на $R$) е релация, съдържаща всички кортежи от $\pi_{\text{останалите атрибути}}(R)$, за които всяка комбинация с кортеж от $S$ съществува в $R$.
\end{definition}

Частното е свързано с декартовото произведение в смисъл, че ако $R \div S = T$, то винаги е в сила включването $T \times S \subseteq R$. Равенството $T \times S = R$ е изпълнено само при специален случай, когато $R$ съдържа всички комбинации между кортежите на $T$ и $S$.

Интуиция: Ако $R$ съдържа $(student\_id, sport)$, а $S$ съдържа $(sport)$, то $R \div S$ дава студентите, които практикуват всички спортове в $S$.

\textit{Пример:}

\begin{center}
\begin{tabular}{|c|c|}
\hline
\textbf{StudentSports} & \\
\hline
student\_id & sport \\
\hline
1 & badminton \\
2 & cricket \\
2 & badminton \\
4 & cricket \\
\hline
\end{tabular}
\hspace{1cm}
\begin{tabular}{|c|}
\hline
\textbf{Sports} \\
\hline
sport \\
\hline
badminton \\
cricket \\
\hline
\end{tabular}
\hspace{1cm}
\begin{tabular}{|c|}
\hline
\textbf{StudentSports} $\div$ \textbf{Sports} \\
\hline
student\_id \\
\hline
2 \\
\hline
\end{tabular}
\end{center}

Резултатът съдържа само студент 2, защото той е единственият студент, който практикува всички спортове от релацията Sports.

\end{FlushLeft}

\end{document}
